\section{Modifications to chapter 11}

The following modifications shall be applied to chapter 11 (\emph{Classes}
[\texttt{class}]):

\begin{modification}[Modifications to 11.4.1 grammar rules]
Modify the following grammar production:\\

member-declarator :
\tabto{1em}declarator virt-specifier-seq\textsubscript{opt}
\stdremove{function-contract-specifier-seq\textsubscript{opt}} pure-specifier\textsubscript{opt}\\\
\tabto{1em}declarator requires-clause
\stdremove{function-contract-specifier-seq\textsubscript{opt}}\\\
\tabto{1em}declarator brace-or-equal-initializer\\\
\tabto{1em}identifier\textsubscript{opt} attribute-specifier-seq\textsubscript{opt} : constant-expression brace-or-equal-initializer\textsubscript{opt}\\
\end{modification}

\begin{modification}[Removal of 11.4.1/2]
\stdremove{The optional \emph{function-contract-specifier-seq} (9.4.1) in a
\emph{member-declarator} shall be present only if the declarator
declares a function.}
\end{modification}

\begin{modification}[Modification of 11.4.1/10]
A \emph{complete-class} context of a class (template) is a

\begin{itemize}[label=\textemdash]

\item
function body (9.6.1),

\item
default argument (9.3.4.7),

\item
default template argument (13.2),

\item
\emph{noexcept-specifier} (14.5), \stdadd{or}

\item
\stdremove{\emph{function-contract-specifier} (9.4.1), or}

\item
default member initializer
within the \emph{member-specification} of the class or class template.
\end{itemize}

NOTE 4 A \emph{complete-class} context of a nested class is also a
\emph{complete-class} context of any enclosing class, if the nested
class is defined within the \emph{member-specification} of the enclosing class.
\end{modification}

\begin{modification}[Modification of 11.8.1/7]
\textbf{Note to Editors}: Remove all uses of \texttt{pre} and \texttt{post} in
example 4.
\end{modification}

\begin{modification}[Modification of 11.9.3/16]
Member functions (including virtual member functions, 11.7.3) can be called for an object under construction
or destruction. Similarly, an object under construction or destruction can be
the operand of the \texttt{typeid}
operator (7.6.1.8) or of a \texttt{dynamic\_cast} (7.6.1.7). However, if these operations are performed during
evaluation of
\stdadd{a \emph{ctor-initializer} (or in a function called directly or indirectly from a
\emph{ctor-initializer}) before all the
\emph{mem-initializers} for base classes have completed, the program has
undefined behavior.}

\begin{itemize}[label=\textemdash]

\item
\stdremove{
a \emph{ctor-initializer} (or in a function called directly or indirectly from a
\emph{ctor-initializer}) before all the
\emph{mem-initializers} for base classes have completed,} 

\item
\stdremove{a precondition assertion of a constructor, or}

\item
\stdremove{a postcondition assertion of a destructor (9.4.1),}

\end{itemize}

\stdremove{the program has undefined behavior.}
\end{modification}

\begin{modification}[Modification of 11.9.5/4]
Member functions, including virtual functions (11.7.3), can be called during
construction or destruction (11.9.3).  When a virtual function is called
directly or indirectly from a constructor or from a destructor, including during
the construction or destruction of the class’s non-static data
members\stdremove{, or\
during the evaluation of a postcondition assertion of a constructor or a
precondition assertion} \stdremove{of a destructor (9.4.1)}, and the object to which the call
applies is the object (call it x) under construction or destruction, the
function called is the final overrider in the constructor’s or destructor’s
class and not one overriding it in a more-derived class.  If the virtual
function call uses an explicit class member access (7.6.1.5) and the object
expression refers to the complete object of x or one of that object’s base class
subobjects but not x or one of its base class subobjects, the behavior is
undefined.
\end{modification}
